% ================================================================
% ==== FILE: content/k_levels/k2.tex
% ================================================================

% ==============================
%  Ontology of Continua — Core
%  K-level Module: K2
%  File: content/k_levels/k2.tex
%  Status: FULLY DEFINED
% ==============================



\paragraph{Canonical tuple projection note.}
Local K-level displays in this module are projections of the canonical public tuple \(K=(\Omega,\partial\Omega,A,\Theta,P,J,C,k,M)\). When a display omits \(M\), reorders components, or uses level-indexed shorthand, the omitted embedding context is inherited from the surrounding level discussion rather than introduced as a competing definition.

\subsubsection{\texorpdfstring{$K_2$}{K_2} Übersicht}
\label{sec:k2-overview}

Level $K_2$ ist das erste Kontinuum, das \emph{genuine
zweidimensionale Struktur}. Es entsteht durch den Dimensionsübergang
$\Psi_{1\to2}$, wenn die Einschränkungen von $K_1$ nicht mehr ausreichen
stellen inkompatible Variationen dar (Axiom~15). Die neue Achse $A_2$ ist
orthogonal im strukturellen Sinne und kann nicht als Funktion ausgedrückt werden
von $A_1$.

Dadurch wird $K_2$ zur ersten Ebene, auf der:
\begin{itemize}
    \item Es gibt nichttriviale Zyklen,
    \item Zeit $\tau$ entsteht (über stabile $C$-Zyklen),
    \item-Gradienten und -Potenziale bilden zweidimensionale Felder,
    Es treten Defekte, Wirbel und Versetzungen auf.
    \item topologische Phänomene werden bedeutungsvoll.
\end{itemize}

$K_2$ ist die minimale Einstellung für klassisches feldtheoretisches Verhalten und
der Übergang vom BKT-Typ.

% ================================================================
\subsubsection{Zustandsraum \texorpdfstring{$\Omega(K_2)$}{\Omega(K_2)}}

Der Zustandsraum von $K_2$ ist:
\[
\Omega(K_2) = C^0(X_1 \times X_2, V),
\]
ein Raum kontinuierlicher Felder über einer 2D-Domäne.

Haupteigenschaften:
\begin{enumerate}
    \item $\Omega(K_2)$ enthält Konfigurationen mit sinnvollen lokalen
          Gradienten in zwei unabhängigen Richtungen.
    \item Defekte (z. B. Wirbel) liegen in $\partial\Omega(K_2)$, es sei denn
          Sie erfüllen die Zulässigkeitsbeschränkungen von $M_2$.
    \item Es gibt nichttriviale Homotopieklassen:
          \[
          \pi_1(\Omega(K_2)) \neq 0.
          \]
\end{enumerate}

Dies ist die erste Ebene, auf der Geometrie (im schwachen strukturellen Sinne)
entsteht.

% ================================================================
\subsubsection{Grenze \texorpdfstring{$\partial\Omega(K_2)$}{\partial\Omega(K_2)}}

$\partial\Omega(K_2)$ enthält:
\begin{itemize}
    \item diskontinuierliche Felder,
    \item Felder, deren Energiedichte divergiert,
    \item-Konfigurationen, die $\Theta_{\mathrm{grad}}$ überschreiten oder
          $\Theta_{\mathrm{defekt}}$,
    \item-Felder, in denen Zyklen zusammenbrechen ($C \to \varnothing$).
\end{itemize}

Ein besonderer Teil der Grenze entspricht topologischen Defekten, deren
Die Kerngröße kollabiert auf Null, was die Zerstörung der Dimension anzeigt
Kohärenz.

% ================================================================
\subsubsection{Achsen \texorpdfstring{$A(K_2)$}{A(K_2)}}

Achsen bestehen aus:
\[
A(K_2) = \{A_1, A_2\},
\]
wo:
\begin{itemize}
    \item $A_1$ wird von $K_1$ geerbt,
    \item $A_2$ ist eine neue unabhängige Strukturachse.
\end{itemize}

Die Unabhängigkeitsbedingung lautet:
\[
\nexists\, h: \mathbb{R} \to \mathbb{R} \text{ mit } A_2 = h \circ A_1.
\]

Diese Unabhängigkeit ist der formale Kern des Dimensionssprungs.

% ================================================================
\subsubsection{Potentiale \texorpdfstring{$P(K_2)$}{P(K_2)}}

Die Menge der Potenziale umfasst jetzt vollständige 2D-Strukturbegriffe:
\[
P(K_2) =
\{ P_{\mathrm{grad}}, P_{\mathrm{curl}}, P_{\mathrm{div}},
   P_{\mathrm{bond}}, P_{\mathrm{smooth}}, P_{\mathrm{top}} \}.
\]

Interpretation:
\begin{itemize}
    \item-Gradienten in zwei Achsen kodieren lokale Spannungen,
    \item Lockenartige Strukturen kodieren Rotationsfreiheitsgrade,
    \item Divergenzpotentiale kodieren Druck- oder Dilatationsspannung,
    \item $P_{\mathrm{top}}$ erfasst die Energie von Wirbeln und Defekten.
\end{itemize}

Diese Potenziale reproduzieren auf abstrakter Ebene das Strukturelle
Vorläufer der Gebiete der Physik.

% ================================================================
\subsubsection{Schwellenwerte \texorpdfstring{$\Theta(K_2)$}{\Theta(K_2)}}

Zu den Schwellenwerten gehören jetzt:
\[
\Theta(K_2) =
\{\Theta_{\mathrm{grad}}, \Theta_{\mathrm{curl}},
  \Theta_{\mathrm{defekt}}, \Theta_{\mathrm{coh}},
  \Theta_{\mathrm{BKT}}\}.
\]

Ihre Bedeutung:
\begin{itemize}
    \item $\Theta_{\mathrm{grad}}$: maximal zulässige Gradientengröße.
    \item $\Theta_{\mathrm{curl}}$: Rotationstoleranz.
    \item $\Theta_{\mathrm{defect}}$: Schwelle für die Stabilität des Defektkerns.
    \item $\Theta_{\mathrm{coh}}$: Kohärenzschwelle für 2D-Ordnung.
    \item $\Theta_{\mathrm{BKT}}$: Schwelle für das Erscheinen von
          stabile $C$-Zyklen; äquivalent zu $K_c$ im BKT-Theorem von
          Darstellbarkeit.
\end{itemize}

Das Überqueren von $\Theta_{\mathrm{BKT}}$ ergibt die entstehende Zeitachse.

% ================================================================
\subsubsection{Flüsse \texorpdfstring{$J(K_2)$}{J(K_2)}}

Zu den Flüssen gehören:
\[
J(K_2) = \{\partial_{x_1} f,\; \partial_{x_2} f,\;
\text{2D-Verformungen},\; \text{Wirbeldrift},\; \text{Phasendrehungen}\}.
\]

Hier erscheinen:
\begin{itemize}
    \item Divergenz und Curlflüsse,
    \item diffusionsähnliches Gleichgewicht,
    \item Dynamik der Fehlervernichtung und Fehlerlösung.
\end{itemize}

Flüsse definieren die dynamische Zulässigkeit der Evolution in $\Omega(K_2)$.

% ================================================================
\subsubsection{Zyklen \texorpdfstring{$C(K_2)$}{C(K_2)}}

$K_2$ ist die erste Ebene mit nichttrivialen Kreisen:
\[
C(K_2) = \{C_{\mathrm{Wirbel}},\; C_{\mathrm{Phase}},\; C_{\mathrm{Ring}}\}.
\]

Am wichtigsten:
\[
C_{\mathrm{Vortex}}: \oint \nabla \theta \cdot d\ell = 2\pi n.
\]

Gemäß dem Theorem of Representability~5 (BKT):
\begin{itemize}
    \item Für die Kopplung $K < K_c$ sind $C_{\mathrm{Vortex}}$ instabil und
          Verfall.
    \item Für $K > K_c$ werden Zyklen stabil und produzieren global
          Kohärenz.
\end{itemize}

Diese stabilen Zyklen erzeugen den zeitlichen Freiheitsgrad.

% ================================================================
\subsubsection{Zeit \texorpdfstring{$\tau(K_2)$}{\tau(K_2)}}

Die Zeit entsteht zum ersten Mal bei $K_2$.

Aus dem Zeitemergenzsatz:
\[
\tau \text{ existiert genau dann, wenn es einen stabilen Zyklus gibt }
C \text{ mit } \Pi(C) > \Theta_{\mathrm{time}}.
\]

Dabei ist $\Pi(C)$ das Zyklus-Leistungsfunktional.

Also:
\[
\tau(K_2) =
\begin{Fälle}
\text{undefiniert}, & K < K_c, \\
\text{wohldefinierter monotoner Parameter}, & K > K_c.
\end{Fälle}
\]

Dies ist das strukturelle Analogon der zeitlichen Ordnung, die sich daraus ergibt
dauerhafte 2D-Kohärenz.

% ================================================================
\subsubsection{Kontinuität \texorpdfstring{$k(K_2)$}{k(K_2)}}

Die allgemeine Formel:
\[
k(K_2) =
\frac{\mu(\Omega(K_2))}{\mu(S(K_2))}
\cdot
\frac{|A(K_2)|}{|A(M_2)|}
\cdot
\frac{\sum \mathrm{Stab}(C)}{\sum \mathrm{MaxStab}(C)}
\cdot
\left(1 - \frac{\sigma(J)}{\sigma_{\max}}\right)
\cdot
\left(1 - \frac{T}{\Theta_{\mathrm{coh}}}\right)_+.
\]

Interpretation:
\begin{itemize}
    Das \item-Wachstum von $\Omega(K_2)$ erhöht $k$,
    \item Erhöhung der Defektdichte senkt $k$,
    \item stabile Zyklen erhöhen $k$,
    \item turbulenzartige Strömungen reduzieren $k$.
\end{itemize}

Dies quantifiziert die strukturelle Lebensfähigkeit des 2D-Kontinuums.

% ================================================================
\subsubsection{Strukturelle Spannung \texorpdfstring{$T(K_2)$}{T(K_2)}}

Spannungsquellen:
\begin{itemize}
    \item große Farbverläufe,
    \item Wirbelkerne,
    \item-Defekt-Anti-Defekt-Interaktionen,
    \item inkompatible Randbedingungen.
\end{itemize}

Ein Kollaps tritt auf, wenn:
\[
T(K_2) > \Theta_{\mathrm{coh}}.
\]

Dies führt zum Kohärenzverlust und zur Zerstörung von $C$-Zyklen.

% ================================================================
\subsubsection{Energie \texorpdfstring{$E(K_2)$}{E(K_2)}}

Die Strukturenergie:
\[
E(K_2) = \int_{X_1 \times X_2}
\mathcal{E}(f, \nabla f, \nabla^2 f)\, d\mu,
\]
wobei $\mathcal{E}$ Folgendes beinhaltet:
\begin{itemize}
    \item Gradientenenergie,
    \item elastische Energie,
    \item Defektkernenergie,
    \item topologische Energie von Wirbeln.
\end{itemize}

Dies ist der strukturelle Prototyp der Energiefunktionale in der Physik.

% ================================================================
\subsubsection{Operatoren auf \texorpdfstring{$K_2$}{K_2} (\texorpdfstring{$\Psi$}{\Psi}, \texorpdfstring{$\Phi$}{\Phi}, \texorpdfstring{$\Lambda$}{\Lambda}, \texorpdfstring{$U$}{U}, \texorpdfstring{$\Chi$}{\Chi})}

\begin{itemize}
    \item $\Psi_{2\to3}$ – Dimensionsübergang zu $K_3$ bei 2D
          Die Struktur wird unzureichend.
    \item $\Phi$ – Metaraum-Evolutionsoperator für $M_2$.
    \item $\Lambda$ – Operator zur Durchsetzung von Einschränkungen, der stabilisiert
          Farbverläufe, Glätte und Fehlergrenzen.
    \item $U$ – Vereinigung lokaler 2D-Regionen zu globalen kohärenten Regionen
          Felder.
    \item $\Chi$ – Verzweigungsoperator, der mehrere $K_2$ ermöglicht
          Mannigfaltigkeiten mit unterschiedlichen Kohärenzstrukturen.
\end{itemize}

% ================================================================
\subsubsection{Prozesse auf \texorpdfstring{$K_2$}{K_2}}

Zu den Prozessen gehören:
\begin{itemize}
    \item Wirbelbildung, Vernichtung und Drift,
    \item Phasenordnung und -unordnung,
    \item Übergänge vom Typ BKT,
    \item Stabilisierung der globalen Kohärenz,
    \item Kollaps unter Defektvermehrung,
    \item Dimensionswachstum auf $K_3$ bei inkompatiblen Variationen
          kann nicht von $\{A_1, A_2\}$ erfasst werden.
\end{itemize}

% ================================================================
\subsubsection{Vorhersagen für \texorpdfstring{$K_2$}{K_2}}

\begin{enumerate}
    \item Vorhandensein einer Kohärenzschwelle $K_c$ (BKT).
    \item Scharfer Übergang zwischen stabilen und instabilen Wirbeln.
    \item Entstehung der Zeit erst, wenn sich die Wirbelzyklen stabilisieren.
    \item Universelle Potenzgesetzkorrelationen unter $K_c$.
    \item Kohärenzverlust unter Defektproliferation.
    \item Notwendigkeit von zwei unabhängigen Achsen für nichttriviale Zyklen.
\end{enumerate}

Diese Vorhersagen stimmen mit einer breiten Klasse physikalischer Systeme überein (XY-Modell,
supraflüssige Filme, dünne magnetische Schichten).

% ================================================================
\subsubsection{Experimente für \texorpdfstring{$K_2$}{K_2}}

Kanonische experimentelle Signaturen:
\begin{itemize}
    \item Messung der Wirbellösungstemperatur in 2D-Supraflüssigkeiten,
    \item Nachweis von BKT-ähnlichen Übergängen in atomar dünnen Filmen,
    \item Beobachtung des Kohärenzkollapses bei hoher Defektdichte,
    \item Phasenordnungskinetik in 2D-Materialien,
    \item Universeller Steifigkeitssprung am BKT-Übergang.
\end{itemize}

Indirekte Tests erscheinen in:
\begin{itemize}
    \item Graphen-Defektdynamik,
    \item 2D-Flüssigkristallfilme,
    \item biologische Membranen mit topologischen Defektmustern.
\end{itemize}

% ================================================================
\subsubsection{Zusammenbruch und Tod von \texorpdfstring{$K_2$}{K_2}}

$K_2$ bricht zusammen, wenn:
\[
\Omega(K_2) = \varnothing,
\]
aufgrund von:
\begin{itemize}
    \item außer Kontrolle geratene Defektvermehrung,
    \item Kohärenzverlust ($T > \Theta_{\mathrm{coh}}$),
    \item Zerstörung aller $C$-Zyklen,
    \item Unfähigkeit, zwei unabhängige Achsen aufrechtzuerhalten.
\end{itemize}

Nach dem Zusammenbruch ist kein Übergang zu $K_3$ möglich.

% ================================================================
\subsubsection{Falschbarkeit von \texorpdfstring{$K_2$}{K_2}}

Falsifizierbare Vorhersagen:
\begin{enumerate}
    \item Existenz einer scharfen BKT-ähnlichen Schwelle,
    \item Notwendigkeit von Zyklen für die entstehende Zeit,
    \item Korrelationslängenverhalten in der Nähe von $K_c$,
    \item Universeller Steifigkeitssprung,
    \item topologischer Schutz von $C_{\mathrm{vortex}}$ über $K_c$.
\end{enumerate}

Eine empirische Verletzung dieser Merkmale würde der Struktur widersprechen
Modell.

% ================================================================
\subsubsection{Verzweigung / ontologische Position von \texorpdfstring{$K_2$}{K_2}}

$K_2$ ist die Wurzel aller höherdimensionalen Kontinua:
\[
K_0 \to K_1 \to K_2 \to K_3 \to \dots
\]

Eine Verzweigung bei $K_2$ entspricht unterschiedlicher zulässiger 2D-Kohärenz
Strukturen:
\begin{itemize}
    \item wirbelreiche vs. wirbelfreie Phasen,
    \item glatte vs. defektdominierte Mannigfaltigkeiten,
    \item Zweige mit hoher Kohärenz vs. Zweige mit niedriger Kohärenz.
\end{itemize}

Diese Zweige bleiben unterschiedlich, sofern sie nicht durch den Operator $U$ vereinheitlicht werden.

% ================================================================
\subsubsection{Beziehung zu M-Räumen}

$K_2$ ist eingebettet in:
\[
K_2 \subseteq M_2 \subseteq M_1 \subseteq M_0.
\]

$M_2$ definiert:
\begin{itemize}
    \item zulässige 2D-Topologien,
    \item erlaubte Fehlerklassen,
    \item Farbverlaufs- und Lockengrenzen,
    \item Kohärenzbedingungen für Dimensionsstabilität.
\end{itemize}

Der Übergang zu $K_3$ erfordert, dass $M_2$ keine neuen Daten aufnehmen kann
inkompatible Differenzklassen.
